\subsection*{A Non-Trivial $d_3$}

We start by listing a view facts. Our main theorem is: 
\begin{theorem}[Summary of the Specrtal Sequence]
	For each Morse datum $(M,f,g,\mathfrak{A})$ together with a filtration $N_{-1}\subseteq N_0 \subseteq \cdots \subseteq N_{m-1} \subseteq N_m=M$ of conlay pairs, there esists a Spectralsequence $(EM^p_r,dM^p_r)$ where each page is graded in $p$ such that
	\begin{itemize}
		\item The first page is the morse complex, hence the second page is the Morse homology
		\item The spectral sequence becomes stationary and the infinity page is the associeted graded goup of the K-groups meaning:
		\begin{align*}
			\bigoplus_{p\in 2\N\phantom{+1}} E_\infty ^p \text{ is a associated graded group to }K^0(M) \\
			\bigoplus_{p\in 2\N+1} E_\infty ^p \text{ is a associated graded group to }K^1(M)\,.
		\end{align*}
	\end{itemize}
	\end{theorem}

\begin{theorem}[Morse-theoretic Cohomological Künneth Formular]
		If $\Lambda$ is a principle ideal domain, and let $(M,\mathfrak{A},\Lambda)$ and $(M',\mathfrak{A}',\Lambda)$ be two Morse data. Then for each $n$ there is a splitting short exact sequence 
	\begin{multline*}
		0~\to ~
		\bigoplus_i\left(H^i(M,\mathfrak{A},\Lambda)\otimes_\Lambda H^{n-i}(M',\mathfrak{A}',\Lambda)\right) ~\to~
		H^n(M\times M',\mathfrak{A}\times \mathfrak{A}',\Lambda) \\
		~	\to ~
		\bigoplus_i\left(\Tor_{\Lambda}\left(
		H^i(M,\mathfrak{A},\Lambda),H^{n-i+1}(M',\mathfrak{A}',\Lambda)
		\right) \right)~ \to ~0
	\end{multline*}
\end{theorem}
\begin{proof}
    The proof is not so bad and will be done in the final thesis. (We need it for a result wich could also be deduced from the isomorohism between Morse and singular cohomology).
\end{proof}

The next theorem comes from Atiyah himself and im not aware of an proof that does not rely an CW-structures:

\begin{theorem}[Künneth Formular for complex K-Theory]
    Let $X;Y$ be finite CW-complexes, then we have a natural exact sequence
    \begin{center}
        \begin{tikzcd}
            0\arrow[r]K^*(X)\otimes K^*(Y)\arrow[r,"\alpha"]& K^*(X\times Y)\arrow[r,"\beta"]&\Tor(K^*(X),K^*(Y))\to 0\,.
        \end{tikzcd}
    \end{center}
    Here $\otimes$ and $\Tor$ are applied to abelian groups, and $\alpha$ is the natural map inducved by the product in $K^*$. Morover, the sequence is $\Z_2$ graded with $\deg(\alpha)=0$, $\deg(\beta)=1$,ehere $K^p\otimes K^q$ and $\Tor(K^p,K^q)$ are given degree $p+q$ ($p,q\in \Z_2$). 
\end{theorem}
We will also use the known $K$-groups of real projective space. We cannot get them through our Spectral sequence direct, as we only get associated groups.
	\begin{fact}
		In fact $K^0(\R\PP^n)\cong \Z \oplus \Z_{2^{\left\lfloor \frac{n}{2}\right\rfloor}}$ and $K^1(\R\PP^n)\cong \Z$ if $n$ is odd, and $0$ if $n$ is even.
	\end{fact}

\begin{example}[$\R\PP^2\times \R\PP^4$]
    First we list the Cohomology of $\R\PP^2$ and $\R\PP^4$ (This is not to bad to calculate by Morse Theory and can be done by reusing earlyer work, where I looked at the Morse-Theoretic Spectral Sequence of $\R\PP^n$):
	\begin{equation*}
		\begin{array}{c|c c}
			& \R\PP^2 &\R\PP^4\\
			\hline
			H^0&\Z&\Z \\
			H^1&0&0\\
			H^2&\Z_2&\Z_2\\
			H^3&0&0\\
			H^4&0&\Z_2
		\end{array}
	\end{equation*}

    Furthermore, the only non-zero $\Tor$ is $\Tor(\Z_2,\Z_2)$, which appears for $n=3$ and $n=5$. For the tensor groups for any $A$ we have: $\Z\otimes A\cong A$ and $\Z_2\otimes \Z_2 \cong \Z_2$. 

    Hence, we have the resulting isomorphisms: 
 \begin{equation*}
 	H^i(\R\PP^2\times \R\PP^4)\cong \begin{cases}
 		\Z & \text{ if }i=  0\,,\\
	 	0  & \text{ if }i=  1\,,\\
	 	\Z_2\oplus \Z_2& \text{ if }i= 2\,,\\
	 	\Z_2 & \text{ if }i=  3\,,\\
	 	\Z_2\oplus \Z_2 & \text{ if }i=  4\,,\\
	 	\Z_2& \text{ if }i=  5\,,\\
	 	\Z_2& \text{ if }i=  6\,,\\
	 	0& \text{ if }i>6\,.
 	\end{cases}
 \end{equation*}
(This result could also have been derived from the classical universal coefficient theorem in singular homology and its isomorphism to Morse homology)
 Now $d_3$ is a map of those groubs of degree three. Hence, there are two maps, that could be non-trivial:
\begin{align}
    d_3^0: &\Z\phantom{_2\oplus \Z_2}\to \Z_2\\
    d_3^2: &\Z_2\oplus \Z_2 \to \Z_2\\
    d_3^3: &\Z_2\phantom{~\oplus \Z_2}\to \Z_2\,.
\end{align}

Now by the Künneth formular for K-Theory, we get the two exact sequences:
\begin{center}
        	\begin{tikzcd}
		0 \arrow[r]&K^0(\R\PP^2)\otimes K^0(\R\PP^4) \arrow[r]&K^0(\R\PP^2\times \R\PP^4)\arrow[r]&\Tor(K^1(\R\PP^2),K^1(\R\PP^4))\arrow[r]&0\\
			0 \arrow[r]&K^1(\R\PP^2)\otimes K^1(\R\PP^4) \arrow[r]&K^1(\R\PP^2\times \R\PP^4)\arrow[r]&\Tor(K^0(\R\PP^2),K^0(\R\PP^4))\arrow[r]&0
	\end{tikzcd}

	\end{center}In our case the two sequences are of the form:\begin{center}
		\begin{tikzcd}
		0 \arrow[r]&(\Z\oplus \Z_2)\otimes (\Z\oplus \Z_4) \arrow[r]&K^0(\R\PP^2\times \R\PP^4)\arrow[r]&\Tor(0,0)\arrow[r]&0\\
		0 \arrow[r]&0\arrow[r]&K^1(\R\PP^2\times \R\PP^4)\arrow[r]&\Tor(\Z_2,\Z_4)\arrow[r]&0
	\end{tikzcd}
	\end{center}
\end{example}

And thereby \begin{equation*}
		K^0(\R\PP^2\times \R\PP^4)\cong \Z\oplus \Z_2\oplus \Z_2\oplus \Z_4\text{ and }	K^1(\R\PP^2\times \R\PP^4)\cong \Z_2
\end{equation*}
Now assuming that all $d_3$ are trivial(and all higher $d_r$ ). Then, the second page would be the last page, where the sum of all odd entrys is a graded associated group of $K^1(\R\PP^2\times \R\PP^4)\cong \Z_2$. Now, under group extensions, the order of the group is multiplicative. Hence, the group assoziated to $K^1(\R\PP^2\times \R\PP^4)$ needs to have order two giving us a contradiction. I Think there are algebaic ways to show that $d_3$ already is nontrivial, by proofing that $d_5^0=0$ (and thereby $d_3$ would be the only possibility for nintrivial maps), as this is a map on the edge of the spectral sequence.

